%2multibyte Version: 5.50.0.2960 CodePage: 65001
% !TeX root = ./Slides_Lez_11_Programmazione_Lineare.tex
%******************
% Slides Lezione 11 - Programmazione Lineare'
% Progetto MQF - Metodi Quantitativi per la Finanza
% Macro matematiche ufficiali del progetto MQF
% Coerenti con MQF_Notazione.tex
%******************
% Fallback utile per compilazione esterna se qualche titolo resta scritto come \QTR.
% Scientific WorkPlace usa \QTR internamente; in Beamer \QTR{frametitle}{Titolo}
% deve espandersi in \frametitle{Titolo}.
%\usepackage{hyperref}
%**** Cartella delle figure: percorso relativo (Overleaf) e assoluto (SWP 5.5)
%**** Titolo, autori e barra inferiore: campo TeX "Frontespizio" dopo \begin{document}

\documentclass[handout]{beamer}%
\usepackage{mathpazo}
\usepackage{multimedia}
\usepackage{amsmath}
\usepackage{amsfonts}
\usepackage{amssymb}
\usepackage{graphicx}%
\setcounter{MaxMatrixCols}{30}
%TCIDATA{OutputFilter=latex2.dll}
%TCIDATA{Version=5.50.0.2960}
%TCIDATA{Codepage=65001}
%TCIDATA{CSTFile=beamer.cst}
%TCIDATA{LastRevised=Saturday, November 14, 2026 10:13:13}
%TCIDATA{<META NAME="GraphicsSave" CONTENT="32">}
%TCIDATA{<META NAME="SaveForMode" CONTENT="2">}
%TCIDATA{BibliographyScheme=Manual}
%TCIDATA{<META NAME="DocumentShell" CONTENT="Other Documents\SW\Slides - Beamer">}
%BeginMSIPreambleData
\providecommand{\U}[1]{\protect\rule{.1in}{.1in}}
%EndMSIPreambleData
\newenvironment{stepenumerate}{\begin{enumerate}[<+->]}{\end{enumerate}}
\newenvironment{stepitemize}{\begin{itemize}[<+->]}{\end{itemize} }
\newenvironment{stepenumeratewithalert}{\begin{enumerate}[<+-| alert@+>]}{\end{enumerate}}
\newenvironment{stepitemizewithalert}{\begin{itemize}[<+-| alert@+>]}{\end{itemize} }
\usetheme{Madrid}
\providecommand{\R}{\mathbb{R}}
\providecommand{\N}{\mathbb{N}}
\providecommand{\E}{\mathbb{E}}
\providecommand{\Prob}{\mathbb{P}}
\providecommand{\Q}{\mathbb{Q}}
\providecommand{\F}{\mathcal{F}}
\providecommand{\B}{\mathcal{B}}
\providecommand{\G}{\mathcal{G}}
\providecommand{\Var}{\operatorname{Var}}
\providecommand{\Cov}{\operatorname{Cov}}
\providecommand{\VaR}{\operatorname{VaR}}
\providecommand{\CVaR}{\operatorname{CVaR}}
\providecommand{\argmin}{\operatorname*{arg\,min}}
\providecommand{\argmax}{\operatorname*{arg\,max}}
\providecommand{\EVPI}{\operatorname{EVPI}}
\providecommand{\VSS}{\operatorname{VSS}}
\providecommand{\QTR}[2]{\csname #1\endcsname{#2}}
\graphicspath{{../graphics/}{E:/Didattica/MQF/Github/CORSO-MQF/graphics/}}
%BeginMSIPreambleData
\ifx\pdfoutput\relax\let\pdfoutput=\undefined\fi
\newcount\msipdfoutput
\ifx\pdfoutput\undefined\else
\ifcase\pdfoutput\else
\msipdfoutput=1
\ifx\paperwidth\undefined\else
\ifdim\paperheight=0pt\relax\else\pdfpageheight\paperheight\fi
\ifdim\paperwidth=0pt\relax\else\pdfpagewidth\paperwidth\fi
\fi\fi\fi
%EndMSIPreambleData
\begin{document}

%****************************************************************
%**** PAGINA DI TITOLO E BARRA INFERIORE (campo TeX "Frontespizio")
%**** Contiene: footline, title, subtitle, author, institute, date,
%**** titlegraphic, maketitle. Resta qui, nel corpo del documento:
%**** nel preambolo SWP 5.5 altera argomenti opzionali e fine riga.
%**** Vedi MQF_Project_Guidelines.md, sezione 10.2.
%****************************************************************%
%TCIMACRO{\TeXButton{Frontespizio}{\setbeamertemplate{footline}{\leavevmode
%\hbox{\begin{beamercolorbox}[wd=.20\paperwidth
%,ht=2.25ex,dp=1ex,center]{author in head/foot}\usebeamerfont
%{author in head/foot}\insertshortauthor\end{beamercolorbox}%
%\begin{beamercolorbox}[wd=.69\paperwidth
%,ht=2.25ex,dp=1ex,center]{title in head/foot}\usebeamerfont
%{title in head/foot}\insertshorttitle\end{beamercolorbox}%
%\begin{beamercolorbox}[wd=.11\paperwidth
%,ht=2.25ex,dp=1ex,right]{date in head/foot}\usebeamerfont{date in head/foot}%
%\insertshortdate{}\hspace*{2em}\insertframenumber{} / \inserttotalframenumber
%\hspace*{2ex}\end{beamercolorbox}}\vskip0pt}\title
%[Lez. 11 - Ottimizzazione e programmazione lineare]{Lezione 11 - Ottimizzazione e programmazione lineare}%
%\subtitle{Metodi Quantitativi per la Finanza}\author
%[P. Falbo, C. Pelizzari]{{\LARGE Paolo Falbo, Cristian Pelizzari}}%
%\institute[]{{\large Dipartimento di Economia e Management}}\date
%[]{}\titlegraphic{\includegraphics[height=1.25cm]{MQF_logo.png}\hspace
%{0.8cm}\includegraphics[height=1.25cm]{logo-Unibs-esteso.png}}\maketitle}}%
%BeginExpansion
\setbeamertemplate{footline}{\leavevmode\hbox{\begin{beamercolorbox}%
[wd=.20\paperwidth,ht=2.25ex,dp=1ex,center]{author in head/foot}%
\usebeamerfont{author in head/foot}\insertshortauthor\end{beamercolorbox}%
\begin{beamercolorbox}[wd=.69\paperwidth
,ht=2.25ex,dp=1ex,center]{title in head/foot}\usebeamerfont
{title in head/foot}\insertshorttitle\end{beamercolorbox}%
\begin{beamercolorbox}[wd=.11\paperwidth
,ht=2.25ex,dp=1ex,right]{date in head/foot}\usebeamerfont{date in head/foot}%
\insertshortdate{}\hspace*{2em}\insertframenumber{} / \inserttotalframenumber
\hspace*{2ex}\end{beamercolorbox}}\vskip0pt}\title
[Lez. 11 - Ottimizzazione e programmazione lineare]{Lezione 11 - Ottimizzazione e programmazione lineare}%
\subtitle{Metodi Quantitativi per la Finanza}\author
[P. Falbo, C. Pelizzari]{{\LARGE Paolo Falbo, Cristian Pelizzari}}%
\institute[]{{\large Dipartimento di Economia e Management}}\date
[]{}\titlegraphic{\includegraphics[height=1.25cm]{MQF_logo.png}\hspace
{0.8cm}\includegraphics[height=1.25cm]{logo-Unibs-esteso.png}}\maketitle
%EndExpansion
%****************************************************************
%**** FINE PAGINA DI TITOLO
%****************************************************************


%**** Compatibilita' Overleaf: ignora natwidth/natheight scritti da SWP
%**** negli \includegraphics (non supportati dal pdftex.def attuale).%
%TCIMACRO{\TeXButton{Grafica natwidth}{\makeatletter\let\KV@Gin@natwidth
%\@gobble\let\KV@Gin@natheight\@gobble\makeatother}}%
%BeginExpansion
\makeatletter\let\KV@Gin@natwidth\@gobble\let\KV@Gin@natheight\@gobble
\makeatother
%EndExpansion


%****************************************


\section{Apertura}

\subsection{Dalla misurazione alla decisione}%

%TCIMACRO{\TeXButton{BeginFrame}{\begin{frame}}}%
%BeginExpansion
\begin{frame}%
%EndExpansion
%

%TCIMACRO{\QTR{frametitle}{Dalla misurazione del rischio alla decisione}}%
%BeginExpansion
\frametitle{Dalla misurazione del rischio alla decisione}%
%EndExpansion
%

%TCIMACRO{\TeXButton{Transparency}{\setbeamercovered{transparent=20}}}%
%BeginExpansion
\setbeamercovered{transparent=20}%
%EndExpansion


Nei Capitoli~8-10, le catene di Markov hanno descritto l'evoluzione degli
stati creditizi; associando a ciascuno stato una perdita monetaria, sono state
costruite distribuzioni di perdita e misure di rischio.

\medskip

{\small
\[%
\begin{array}
[c]{rcl}%
\mbox{Cap. 8-10} & : & \mbox{stati futuri}\rightarrow
\mbox{probabilit\`a}\rightarrow\mbox{perdite}\rightarrow
\mbox{misure di rischio}\\[8pt]%
\mbox{Cap. 11} & : & \mbox{variabili decisionali}\rightarrow
\mbox{vincoli}\rightarrow\mbox{decisioni ammissibili}\\
&  & \rightarrow\mbox{scelta ottima}\rightarrow\mbox{sensitivit\`a}\rightarrow
\mbox{valori marginali dei vincoli}
\end{array}
\]
}

\medskip

\begin{stepitemize}
\item La prima catena dice quali eventi avversi possono verificarsi, con quale
probabilit\`{a} e quanto pesa la parte pi\`{u} sfavorevole della distribuzione.

\item Non stabilisce per\`{o} \textbf{quale decisione} debba essere adottata.

\item Un intermediario deve impiegare le risorse \emph{prima} che gli eventi
futuri siano osservati: scegliere le posizioni da assumere, quante risorse
mantenere liquide, quale redditivit\`{a} perseguire.
\end{stepitemize}

%

%TCIMACRO{\TeXButton{Transition: Box Out}{\transboxout}}%
%BeginExpansion
\transboxout
%EndExpansion%
%TCIMACRO{\TeXButton{EndFrame}{\end{frame}}}%
%BeginExpansion
\end{frame}%
%EndExpansion


%****************************************
%

%TCIMACRO{\TeXButton{BeginFrame}{\begin{frame}}}%
%BeginExpansion
\begin{frame}%
%EndExpansion
%

%TCIMACRO{\QTR{frametitle}{Percorso della lezione}}%
%BeginExpansion
\frametitle{Percorso della lezione}%
%EndExpansion
%

%TCIMACRO{\TeXButton{Transparency}{\setbeamercovered{transparent=20}}}%
%BeginExpansion
\setbeamercovered{transparent=20}%
%EndExpansion


\begin{stepenumerate}
\item[1.] \textbf{Formulare}: variabili decisionali, parametri, funzione
obiettivo e vincoli di un modello bancario stilizzato.

\item[2.] \textbf{Verificare}: quali composizioni dell'attivo sono
ammissibili, quali vincoli sono attivi e quali conservano un margine.

\item[3.] \textbf{Vedere}: la regione ammissibile, i punti estremi e le rette
di livello in due dimensioni.

\item[4.] \textbf{Leggere l'ottimo}: non soltanto il valore della funzione
obiettivo, ma la composizione della decisione e i vincoli che la determinano.

\item[5.] \textbf{Misurare il costo della prudenza}: di quanto si riduce il
risultato economico quando i requisiti diventano pi\`{u} severi.

\item[6.] \textbf{Interpretare i vincoli al margine}: collegare la
sensitivit\`{a} del valore ottimo ai prezzi ombra e leggere il costo marginale
dei requisiti che condizionano la soluzione.
\end{stepenumerate}

\medskip

\begin{center}
{\small \emph{Chiave di lettura: il vincolo non \`{e} un dettaglio tecnico
della formulazione, \`{e} la decisione.}}
\end{center}

%

%TCIMACRO{\TeXButton{Transition: Box Out}{\transboxout}}%
%BeginExpansion
\transboxout
%EndExpansion%
%TCIMACRO{\TeXButton{EndFrame}{\end{frame}}}%
%BeginExpansion
\end{frame}%
%EndExpansion


\section{Il caso Silicon Valley Bank}

%****************************************


\subsection{Raccolta e attivo}%

%TCIMACRO{\TeXButton{BeginFrame}{\begin{frame}}}%
%BeginExpansion
\begin{frame}%
%EndExpansion
%

%TCIMACRO{\QTR{frametitle}{Silicon Valley Bank: la raccolta e l'attivo}}%
%BeginExpansion
\frametitle{Silicon Valley Bank: la raccolta e l'attivo}%
%EndExpansion
%

%TCIMACRO{\TeXButton{Transparency}{\setbeamercovered{transparent=20}}}%
%BeginExpansion
\setbeamercovered{transparent=20}%
%EndExpansion


\begin{stepitemize}


\item \textbf{Crescita.} Tra il 2019 e il 2021 le attivit\`{a} del gruppo
passano da circa $71$ a oltre $211$ miliardi di dollari: un volume molto
elevato di nuove risorse da investire in un periodo breve.

\item \textbf{Concentrazione della raccolta.} I depositanti appartengono in
larga parte agli stessi settori e possono manifestare esigenze di
liquidit\`{a} correlate. A fine $2022$ circa il $94\%$ dei depositi non \`{e}
coperto dall'assicurazione federale.

\item \textbf{Impieghi.} Una parte consistente delle risorse \`{e} investita
in titoli pubblici e \emph{mortgage-backed securities} a lunga scadenza:
rischio di credito contenuto, ma elevata sensibilit\`{a} alle variazioni dei tassi.

\item \textbf{Lo scarto del $2022$.} L'aumento dei tassi riduce il valore di
mercato dei titoli lunghi; il rallentamento dei finanziamenti alle imprese
tecnologiche induce i clienti a utilizzare i depositi.
\end{stepitemize}

\medskip

\begin{center}
{\small Rischio di credito, rischio di tasso e rischio di liquidit\`{a} sono
\textbf{dimensioni distinte}: un titolo pu\`{o} essere sicuro rispetto al
rimborso del capitale e perdere valore quando i tassi aumentano.}
\end{center}

%

%TCIMACRO{\TeXButton{Transition: Box Out}{\transboxout}}%
%BeginExpansion
\transboxout
%EndExpansion%
%TCIMACRO{\TeXButton{EndFrame}{\end{frame}}}%
%BeginExpansion
\end{frame}%
%EndExpansion


%***********************************************************


%****************************************


\subsection{Il meccanismo della fragilit\`{a}}%

%TCIMACRO{\TeXButton{BeginFrame}{\begin{frame}}}%
%BeginExpansion
\begin{frame}%
%EndExpansion
%

%TCIMACRO{\QTR{frametitle}{Il meccanismo della fragilit\`{a}}}%
%BeginExpansion
\frametitle{Il meccanismo della fragilit\`{a}}%
%EndExpansion
%

%TCIMACRO{\TeXButton{Transparency}{\setbeamercovered{transparent=20}}}%
%BeginExpansion
\setbeamercovered{transparent=20}%
%EndExpansion


\begin{center}%
%TCIMACRO{\FRAME{itbpF}{4.1494in}{1.9726in}{0in}{}{}%
%{cap11_svb_meccanismo_fragilita.png}{\special{ language "Scientific Word";
%type "GRAPHIC";  maintain-aspect-ratio TRUE;  display "USEDEF";
%valid_file "F";  width 4.1494in;  height 1.9726in;  depth 0in;
%original-width 7.53in;  original-height 3.5533in;  cropleft "0";
%croptop "1";  cropright "1";  cropbottom "0";
%filename '../graphics/Cap11_SVB_meccanismo_fragilita.png';file-properties "XNPEU";}%
%} }%
%BeginExpansion
{\includegraphics[
natheight=3.553300in,
natwidth=7.530000in,
height=1.9726in,
width=4.1494in
]%
{../graphics/Cap11_SVB_meccanismo_fragilita.png}%
}
%EndExpansion



\end{center}

\begin{stepitemize}
\item La fragilit\`{a} non dipende da una singola classe di attivo, ma dalla
\textbf{relazione} tra struttura dell'attivo e struttura della raccolta.

\item Nel marzo $2023$ l'annuncio della vendita di una parte del portafoglio
(in perdita) rende visibili le difficolt\`{a}: il $9$ marzo vengono richiesti
prelievi per oltre $40$ miliardi, il giorno successivo la banca viene chiusa.
\end{stepitemize}

%

%TCIMACRO{\TeXButton{Transition: Box Out}{\transboxout}}%
%BeginExpansion
\transboxout
%EndExpansion%
%TCIMACRO{\TeXButton{EndFrame}{\end{frame}}}%
%BeginExpansion
\end{frame}%
%EndExpansion


%***********************************************************


%****************************************


\subsection{Il problema guida}%

%TCIMACRO{\TeXButton{BeginFrame}{\begin{frame}}}%
%BeginExpansion
\begin{frame}%
%EndExpansion
%

%TCIMACRO{\QTR{frametitle}{Quattro elementi, un problema decisionale}}%
%BeginExpansion
\frametitle{Quattro elementi, un problema decisionale}%
%EndExpansion
%

%TCIMACRO{\TeXButton{Transparency}{\setbeamercovered{transparent=20}}}%
%BeginExpansion
\setbeamercovered{transparent=20}%
%EndExpansion


{\small
\[%
\begin{array}
[c]{c}%
\mbox{rapida crescita delle risorse da investire}\\[3pt]%
\mbox{concentrazione e instabilit\`{a} potenziale della raccolta}\\[3pt]%
\mbox{investimento in attivit\`{a} a lunga scadenza sensibili ai tassi}\\[3pt]%
\mbox{difficolt\`{a} di liquidazione dell'attivo in presenza di deflussi}
\end{array}
\]
}

\medskip

Considerati separatamente, questi elementi non identificano una decisione:

\begin{stepitemize}
\item pi\`{u} attivit\`{a} liquide riducono la vulnerabilit\`{a} ai deflussi,
ma tendono a ridurre la redditivit\`{a};

\item pi\`{u} titoli a lunga scadenza accrescono il rendimento corrente, ma
aumentano la sensibilit\`{a} ai tassi;

\item un limite pi\`{u} severo alle perdite migliora la resistenza allo
scenario avverso, ma restringe gli impieghi disponibili.
\end{stepitemize}

\medskip

\begin{center}
{\small \emph{Prima che una crisi si manifesti, come pu\`{o} una banca
scegliere la composizione del proprio attivo perseguendo la redditivit\`{a},
ma rispettando simultaneamente requisiti di liquidit\`{a}, limiti alle perdite
di valore e vincoli di concentrazione? Quale costo economico comporta rendere
il bilancio pi\`{u} resiliente?}}
\end{center}

%

%TCIMACRO{\TeXButton{Transition: Box Out}{\transboxout}}%
%BeginExpansion
\transboxout
%EndExpansion%
%TCIMACRO{\TeXButton{EndFrame}{\end{frame}}}%
%BeginExpansion
\end{frame}%
%EndExpansion


%****************************************


\section{Variabili, parametri, risultati}

\subsection{Le variabili decisionali}%

%TCIMACRO{\TeXButton{BeginFrame}{\begin{frame}}}%
%BeginExpansion
\begin{frame}%
%EndExpansion
%

%TCIMACRO{\QTR{frametitle}{Che cosa la banca pu\`{o} decidere}}%
%BeginExpansion
\frametitle{Che cosa la banca pu\`{o} decidere}%
%EndExpansion
%

%TCIMACRO{\TeXButton{Transparency}{\setbeamercovered{transparent=20}}}%
%BeginExpansion
\setbeamercovered{transparent=20}%
%EndExpansion


Tassi, andamento del settore tecnologico e richieste di rimborso non erano
controllabili. La \textbf{composizione dell'attivo} dipendeva invece da
decisioni gestionali. Consideriamo quattro classi aggregate:

\medskip

{\small
\[%
\begin{array}
[c]{ll}\hline
x_{1} & \mbox{Cassa e riserve: disponibili subito, redditivit\`{a} minima}\\
x_{2} &
\mbox{Titoli a breve: monetizzabili in tempi brevi, poco sensibili ai
tassi}\\
x_{3} &
\mbox{Titoli a lunga: pi\`{u} redditizi, valore sensibile a un rialzo
dei tassi}\\
x_{4} & \mbox{Prestiti e impieghi meno liquidi: contributo elevato, non
negoziabili}\\\hline
\end{array}
\]
}

\medskip

\begin{stepitemize}
\item Le $x_{j}$ sono le \textbf{variabili decisionali}: non sono dati
osservati e non sono risultati prodotti dal modello.

\item Poich\'{e} rappresentano importi investiti, $x_{j}\geq0$: la condizione
esclude vendite allo scoperto e posizioni negative.

\item Una politica di allocazione \`{e} un vettore $x=(x_{1},x_{2},x_{3}%
,x_{4})^{\prime}$; vettori diversi sono politiche diverse.
\end{stepitemize}

%

%TCIMACRO{\TeXButton{Transition: Box Out}{\transboxout}}%
%BeginExpansion
\transboxout
%EndExpansion%
%TCIMACRO{\TeXButton{EndFrame}{\end{frame}}}%
%BeginExpansion
\end{frame}%
%EndExpansion


%****************************************


\subsection{I parametri}%

%TCIMACRO{\TeXButton{BeginFrame}{\begin{frame}}}%
%BeginExpansion
\begin{frame}%
%EndExpansion
%

%TCIMACRO{\QTR{frametitle}{I dati che servono per valutare una composizione}}%
%BeginExpansion
\frametitle{I dati che servono per valutare una composizione}%
%EndExpansion
%

%TCIMACRO{\TeXButton{Transparency}{\setbeamercovered{transparent=20}}}%
%BeginExpansion
\setbeamercovered{transparent=20}%
%EndExpansion


Per ogni classe di attivo $j$ introduciamo tre coefficienti:

{\small
\[%
\begin{array}
[c]{lll}\hline
c_{j} & \mbox{contributo economico unitario} & \displaystyle\sum_{j=1}%
^{4}c_{j}x_{j}\\[8pt]%
q_{j} & \mbox{quota liquidabile entro l'orizzonte} & \displaystyle\sum
_{j=1}^{4}q_{j}x_{j}\\[8pt]%
\ell_{j} & \mbox{perdita unitaria nello scenario di stress} &
\displaystyle\sum_{j=1}^{4}\ell_{j}x_{j}\\[4pt]\hline
\end{array}
\]
}

\medskip

\begin{stepitemize}
\item $c_{j}$ non \`{e} necessariamente il solo tasso contrattuale: pu\`{o}
essere un margine economico atteso al netto dei costi attribuibili.

\item $q_{j}$ \textbf{non misura la qualit\`{a} creditizia}. Un titolo pu\`{o}
avere rischio di insolvenza minimo e non essere utile per fronteggiare
pagamenti concentrati in pochi giorni. Il coefficiente ha senso solo rispetto
a un orizzonte dichiarato.

\item I coefficienti sono \textbf{parametri}: non vengono scelti dalla
procedura di ottimizzazione, ma fissati o stimati prima della soluzione.
\end{stepitemize}

%

%TCIMACRO{\TeXButton{Transition: Box Out}{\transboxout}}%
%BeginExpansion
\transboxout
%EndExpansion%
%TCIMACRO{\TeXButton{EndFrame}{\end{frame}}}%
%BeginExpansion
\end{frame}%
%EndExpansion


%****************************************


\subsection{Decisioni, dati, risultati, requisiti}%

%TCIMACRO{\TeXButton{BeginFrame}{\begin{frame}}}%
%BeginExpansion
\begin{frame}%
%EndExpansion
%

%TCIMACRO{\QTR{frametitle}{Quattro componenti da non confondere}}%
%BeginExpansion
\frametitle{Quattro componenti da non confondere}%
%EndExpansion
%

%TCIMACRO{\TeXButton{Transparency}{\setbeamercovered{transparent=20}}}%
%BeginExpansion
\setbeamercovered{transparent=20}%
%EndExpansion


Le tre somme della slide precedente non sono dati: sono \textbf{risultati}
della decisione. Una volta scelto $x$, possono essere calcolate.

\medskip

\begin{center}
{\small
\begin{tabular}
[c]{ll}\hline
\textbf{Variabili decisionali} & che cosa pu\`{o} essere scelto: $x_{1}%
,\ldots,x_{4}$\\[3pt]%
\textbf{Parametri} & fissati prima della soluzione: $c_{j}$, $q_{j}$,
$\ell_{j}$\\[3pt]%
\textbf{Risultati} & che cosa discende dalla scelta: redditivit\`{a},\\
& liquidit\`{a}, perdita sotto stress\\[3pt]%
\textbf{Requisiti} & che cosa rende accettabile una decisione\\\hline
\end{tabular}
}
\end{center}

\medskip

\begin{stepitemize}
\item La banca pu\`{o} perseguire un risultato economico elevato, ma non
pu\`{o} ignorare la necessit\`{a} di mantenere liquidit\`{a} sufficiente e di
contenere le perdite potenziali.

\item Il modello dovr\`{a} quindi stabilire \emph{quale} risultato migliorare
e \emph{quali} livelli minimi o massimi rispettare.
\end{stepitemize}

\medskip

\begin{center}
{\small La programmazione lineare organizza queste quattro componenti in un
unico modello: \textbf{variabili}, \textbf{funzione obiettivo},
\textbf{vincoli}.}
\end{center}

%

%TCIMACRO{\TeXButton{Transition: Box Out}{\transboxout}}%
%BeginExpansion
\transboxout
%EndExpansion%
%TCIMACRO{\TeXButton{EndFrame}{\end{frame}}}%
%BeginExpansion
\end{frame}%
%EndExpansion


%*********************************


\section{Costruzione del modello}

%****************************************


\subsection{Bilancio e funzione obiettivo}%

%TCIMACRO{\TeXButton{BeginFrame}{\begin{frame}}}%
%BeginExpansion
\begin{frame}%
%EndExpansion
%

%TCIMACRO{\QTR{frametitle}{Il vincolo di bilancio e la funzione obiettivo}}%
%BeginExpansion
\frametitle{Il vincolo di bilancio e la funzione obiettivo}%
%EndExpansion
%

%TCIMACRO{\TeXButton{Transparency}{\setbeamercovered{transparent=20}}}%
%BeginExpansion
\setbeamercovered{transparent=20}%
%EndExpansion


Sia $A_{0}$ l'ammontare complessivo delle risorse da allocare. L'intero attivo
considerato deve essere distribuito tra le quattro classi:
\[
x_{1}+x_{2}+x_{3}+x_{4}=A_{0}.
\]


\begin{stepitemize}
\item Il vincolo non dice \emph{come} impiegare le risorse: dice soltanto che
la somma degli impieghi coincide con le risorse disponibili.

\item Le risorse non investite in titoli o prestiti sono rappresentate dalla
variabile relativa alla cassa.
\end{stepitemize}

\medskip

Se la banca intende massimizzare il risultato economico,
\[
\max\; c_{1}x_{1}+c_{2}x_{2}+c_{3}x_{3}+c_{4}x_{4}.
\]


\medskip

\begin{center}
{\small Perseguire la redditivit\`{a} \textbf{non} deve portare ad ignorare il
rischio: redditivit\`{a}, liquidit\`{a} e perdite non vengono sommate in un
unico indicatore. La prima \`{e} il risultato da migliorare, le altre
delimitano ci\`{o} che \`{e} accettabile.}
\end{center}

%

%TCIMACRO{\TeXButton{Transition: Box Out}{\transboxout}}%
%BeginExpansion
\transboxout
%EndExpansion%
%TCIMACRO{\TeXButton{EndFrame}{\end{frame}}}%
%BeginExpansion
\end{frame}%
%EndExpansion


%***********************************************************


%****************************************


\subsection{Il requisito di liquidit\`{a}}%

%TCIMACRO{\TeXButton{BeginFrame}{\begin{frame}}}%
%BeginExpansion
\begin{frame}%
%EndExpansion
%

%TCIMACRO{\QTR{frametitle}{Il requisito minimo di liquidit\`{a}}}%
%BeginExpansion
\frametitle{Il requisito minimo di liquidit\`{a}}%
%EndExpansion
%

%TCIMACRO{\TeXButton{Transparency}{\setbeamercovered{transparent=20}}}%
%BeginExpansion
\setbeamercovered{transparent=20}%
%EndExpansion


La liquidit\`{a} disponibile dalla composizione $x$ entro l'orizzonte
considerato \`{e}
\[
q_{1}x_{1}+q_{2}x_{2}+q_{3}x_{3}+q_{4}x_{4}.
\]


Se $Q_{\min}$ \`{e} il livello minimo necessario per fronteggiare i deflussi
previsti nello scenario considerato,
\[
q_{1}x_{1}+q_{2}x_{2}+q_{3}x_{3}+q_{4}x_{4}\geq Q_{\min}.
\]


\medskip

\begin{stepitemize}
\item Il vincolo esclude le composizioni che, \textbf{pur essendo redditizie},
non rendono disponibile un ammontare sufficiente di risorse nei tempi richiesti.

\item $Q_{\min}$ non \`{e} prodotto dal modello: \`{e} una scelta prudenziale
o gestionale \textbf{prefissata}.

\item Un suo aumento restringe l'insieme delle composizioni ammissibili e
tende a spostare le risorse verso le attivit\`{a} maggiormente liquidabili.
\end{stepitemize}

\medskip

\begin{center}
{\small \`{E} il vincolo su cui torneremo: nella soluzione ottima
risulter\`{a} \textbf{attivo}, e alla fine della lezione ne misureremo il
costo.}
\end{center}

%

%TCIMACRO{\TeXButton{Transition: Box Out}{\transboxout}}%
%BeginExpansion
\transboxout
%EndExpansion%
%TCIMACRO{\TeXButton{EndFrame}{\end{frame}}}%
%BeginExpansion
\end{frame}%
%EndExpansion


%***********************************************************


%****************************************


\subsection{Perdite, concentrazione, non negativit\`{a}}%

%TCIMACRO{\TeXButton{BeginFrame}{\begin{frame}}}%
%BeginExpansion
\begin{frame}%
%EndExpansion
%

%TCIMACRO{\QTR{frametitle}{Perdite sotto stress, concentrazione, non
%negativit\`{a}}}%
%BeginExpansion
\frametitle{Perdite sotto stress, concentrazione, non negativit\`{a}}%
%EndExpansion
%

%TCIMACRO{\TeXButton{Transparency}{\setbeamercovered{transparent=20}}}%
%BeginExpansion
\setbeamercovered{transparent=20}%
%EndExpansion


Se $K$ \`{e} la perdita massima compatibile con la capacit\`{a} patrimoniale,
\[
\ell_{1}x_{1}+\ell_{2}x_{2}+\ell_{3}x_{3}+\ell_{4}x_{4}\leq K.
\]


Per limitare l'esposizione verso i titoli a lunga scadenza,
\[
x_{3}\leq\alpha A_{0}, \qquad0\leq\alpha\leq1.
\]


\smallskip

Infine, poich\'{e} le variabili rappresentano importi investiti,
\[
x_{j}\geq0, \qquad j=1,\ldots,4.
\]


\medskip

\begin{stepitemize}
\item I due limiti controllano vulnerabilit\`{a} \textbf{diverse}: il primo
l'effetto economico dello scenario di stress, il secondo la dipendenza da una
singola classe di attivo \emph{anche quando la perdita stimata resta entro}
$K$.

\item Una composizione pu\`{o} rispettare il requisito di liquidit\`{a} ed
essere comunque inammissibile perch\'{e} troppo esposta alle perdite.
\end{stepitemize}

%

%TCIMACRO{\TeXButton{Transition: Box Out}{\transboxout}}%
%BeginExpansion
\transboxout
%EndExpansion%
%TCIMACRO{\TeXButton{EndFrame}{\end{frame}}}%
%BeginExpansion
\end{frame}%
%EndExpansion


%***********************************************************


%****************************************


\subsection{Il modello SVB stilizzato}%

%TCIMACRO{\TeXButton{BeginFrame}{\begin{frame}}}%
%BeginExpansion
\begin{frame}%
%EndExpansion
%

%TCIMACRO{\QTR{frametitle}{Il modello SVB stilizzato}}%
%BeginExpansion
\frametitle{Il modello SVB stilizzato}%
%EndExpansion
%

%TCIMACRO{\TeXButton{Transparency}{\setbeamercovered{transparent=20}}}%
%BeginExpansion
\setbeamercovered{transparent=20}%
%EndExpansion


{\small
\[%
\begin{array}
[c]{ll}%
\max & c_{1}x_{1}+c_{2}x_{2}+c_{3}x_{3}+c_{4}x_{4}\\[4pt]%
\text{soggetto a} & x_{1}+x_{2}+x_{3}+x_{4}=A_{0},\\[3pt]
& q_{1}x_{1}+q_{2}x_{2}+q_{3}x_{3}+q_{4}x_{4}\geq Q_{\min},\\[3pt]
& \ell_{1}x_{1}+\ell_{2}x_{2}+\ell_{3}x_{3}+\ell_{4}x_{4}\leq K,\\[3pt]
& x_{3}\leq\alpha A_{0},\\[3pt]
& x_{j}\geq0,\qquad j=1,\ldots,4.
\end{array}
\]
}

\medskip

\begin{center}
{\small
\begin{tabular}
[c]{lccc}\hline
\textbf{Classe di attivo} & $c_{j}$ & $q_{j}$ & $\ell_{j}$\\\hline
$x_{1}$ Cassa e riserve & $0.010$ & $1.00$ & $0.00$\\
$x_{2}$ Titoli a breve scadenza & $0.025$ & $0.90$ & $0.02$\\
$x_{3}$ Titoli a lunga scadenza & $0.040$ & $0.50$ & $0.15$\\
$x_{4}$ Prestiti e impieghi meno liquidi & $0.050$ & $0.20$ & $0.08$\\\hline
\end{tabular}
}
\end{center}

\smallskip

\begin{center}
{\small $A_{0}=100$, \quad$Q_{\min}=55$, \quad$K=7$, \quad$\alpha=0.30$.
\smallskip}
\end{center}

{\small Il modello non cerca la composizione pi\`{u} redditizia in assoluto,
ma la pi\`{u} redditizia \textbf{tra quelle che rispettano simultaneamente
tutti i requisiti}.}%

%TCIMACRO{\TeXButton{Transition: Box Out}{\transboxout}}%
%BeginExpansion
\transboxout
%EndExpansion%
%TCIMACRO{\TeXButton{EndFrame}{\end{frame}}}%
%BeginExpansion
\end{frame}%
%EndExpansion


%*****************************


\section{Linearit\`{a} e forma generale}

%****************************************


\subsection{Che cosa rende lineare il modello}%

%TCIMACRO{\TeXButton{BeginFrame}{\begin{frame}}}%
%BeginExpansion
\begin{frame}%
%EndExpansion
%

%TCIMACRO{\QTR{frametitle}{Che cosa rende lineare il modello}}%
%BeginExpansion
\frametitle{Che cosa rende lineare il modello}%
%EndExpansion
%

%TCIMACRO{\TeXButton{Transparency}{\setbeamercovered{transparent=20}}}%
%BeginExpansion
\setbeamercovered{transparent=20}%
%EndExpansion


Ogni variabile compare moltiplicata per un coefficiente noto e i contributi
vengono sommati. La linearit\`{a} non dipende dal fatto che il problema sia
semplice o contenga poche variabili, ma dal \textbf{modo} in cui le decisioni
entrano nella funzione obiettivo e nei vincoli.

\medskip

\begin{center}
{\small
\begin{tabular}
[c]{ll}\hline
\textbf{Proporzionalit\`{a}} & il contributo \`{e} proporzionale al livello:
$c_{j}x_{j}$, $q_{j}x_{j}$, $\ell_{j}x_{j}$\\[3pt]%
\textbf{Additivit\`{a}} & il risultato \`{e} la somma dei contributi, senza
interazioni\\[3pt]%
\textbf{Divisibilit\`a} & le variabili possono assumere valori
frazionari\\[3pt]%
\textbf{Coefficienti} & $c_{j}$, $q_{j}$, $\ell_{j}$ e i livelli $A_{0}$,
$Q_{\min}$, $K$, $\alpha$\\
\quad\textbf{come parametri} & sono fissati prima della soluzione\\\hline
\end{tabular}
}
\end{center}

\medskip

\begin{stepitemize}
\item Nel caso SVB: rendimento, quota liquidabile e perdita unitaria
attribuiti a una classe \textbf{non cambiano} al variare dell'importo investito.

\item La divisibilit\`{a} \`{e} naturale per importi monetari; non lo sarebbe
per decisioni indivisibili, che richiederebbero variabili intere.
\end{stepitemize}

%

%TCIMACRO{\TeXButton{Transition: Box Out}{\transboxout}}%
%BeginExpansion
\transboxout
%EndExpansion%
%TCIMACRO{\TeXButton{EndFrame}{\end{frame}}}%
%BeginExpansion
\end{frame}%
%EndExpansion


%***********************************************************


%****************************************


\subsection{Che cosa la linearit\`{a} esclude}%

%TCIMACRO{\TeXButton{BeginFrame}{\begin{frame}}}%
%BeginExpansion
\begin{frame}%
%EndExpansion
%

%TCIMACRO{\QTR{frametitle}{Che cosa la linearit\`{a} esclude}}%
%BeginExpansion
\frametitle{Che cosa la linearit\`{a} esclude}%
%EndExpansion
%

%TCIMACRO{\TeXButton{Transparency}{\setbeamercovered{transparent=20}}}%
%BeginExpansion
\setbeamercovered{transparent=20}%
%EndExpansion


Non compaiono potenze delle variabili, prodotti tra variabili, coefficienti
dipendenti dalle quantit\`{a} investite. Non sarebbero lineari
\[
x_{3}^{2}, \qquad x_{2}x_{3}, \qquad c_{3}(x_{3})\,x_{3}.
\]


\medskip

\begin{stepitemize}
\item \textbf{Nessuna amplificazione.} Il modello non considera che
un'esposizione simultanea verso titoli a lunga e prestiti poco liquidi possa
produrre una vulnerabilit\`{a} superiore alla somma delle due: un termine
$\gamma x_{3}x_{4}$ renderebbe il modello non lineare.

\item \textbf{No fire sale.} La vendita di una grande posizione pu\`{o}
deprimere il prezzo: la perdita unitaria non resta costante al crescere
dell'importo liquidato. Servirebbe $L_{j}=L_{j}(x_{j})$.

\item \textbf{L'incertezza non scompare}, viene incorporata nei coefficienti.
Il modello risponde a una domanda condizionata: qual \`{e} la migliore
allocazione \emph{se} i coefficienti e lo scenario di stress assumono i valori indicati.
\end{stepitemize}

\medskip%

%TCIMACRO{\TeXButton{Transition: Box Out}{\transboxout}}%
%BeginExpansion
\transboxout
%EndExpansion%
%TCIMACRO{\TeXButton{EndFrame}{\end{frame}}}%
%BeginExpansion
\end{frame}%
%EndExpansion


%****************************************


%****************************************


\subsection{Dalla forma scalare alla forma matriciale}%

%TCIMACRO{\TeXButton{BeginFrame}{\begin{frame}}}%
%BeginExpansion
\begin{frame}%
%EndExpansion
%

%TCIMACRO{\QTR{frametitle}{Dal caso SVB alla formulazione generale}}%
%BeginExpansion
\frametitle{Dal caso SVB alla formulazione generale}%
%EndExpansion
%

%TCIMACRO{\TeXButton{Transparency}{\setbeamercovered{transparent=20}}}%
%BeginExpansion
\setbeamercovered{transparent=20}%
%EndExpansion


La struttura non dipende dal numero di classi. Con $n$ variabili e $m$
requisiti, in forma scalare e in forma matriciale, con $A=(a_{ij})$:
\[%
\begin{array}
[c]{ll}%
\max & \displaystyle\sum_{j=1}^{n}c_{j}x_{j}\\[6pt]%
\text{s.a.} & \displaystyle\sum_{j=1}^{n}a_{ij}x_{j}\leq b_{i},\quad
i=1,\ldots,m,\\[4pt]
& x_{j}\geq0,\quad j=1,\ldots,n,
\end{array}
\qquad\Longleftrightarrow\qquad%
\begin{array}
[c]{ll}%
\max\limits_{x} & c^{\prime}x\\[4pt]%
\text{s.a.} & Ax\leq b,\\[2pt]
& x\geq0.
\end{array}
\]


\begin{stepitemize}
\item Ogni \textbf{riga} di $A$ \`e un requisito; ogni \textbf{colonna}
descrive come una variabile contribuisce ai diversi requisiti. Nel modello SVB
la colonna di una classe di attivo ne raccoglie il peso nel bilancio, la
liquidabilit\`a $q_{j}$, la perdita sotto stress $\ell_{j}$ e, per i titoli a
lunga, l'esposizione al limite di concentrazione. La redditivit\`a $c_{j}$ non
sta in $A$, ma nel vettore $c$ della funzione obiettivo.

\item Uguaglianze e vincoli $\geq$ si riconducono a questa forma: ad esempio
$q^{\prime}x\geq Q_{\min}$ equivale a $-q^{\prime}x\leq-Q_{\min}$.
\end{stepitemize}

%

%TCIMACRO{\TeXButton{Transition: Box Out}{\transboxout}}%
%BeginExpansion
\transboxout
%EndExpansion%
%TCIMACRO{\TeXButton{EndFrame}{\end{frame}}}%
%BeginExpansion
\end{frame}%
%EndExpansion


%*****************************


\section{Decisioni ammissibili e diagnosi}

%*****************************


\subsection{Decisioni ammissibili}%

%TCIMACRO{\TeXButton{BeginFrame}{\begin{frame}}}%
%BeginExpansion
\begin{frame}%
%EndExpansion
%

%TCIMACRO{\QTR{frametitle}{Quali composizioni possono essere considerate}}%
%BeginExpansion
\frametitle{Quali composizioni possono essere considerate}%
%EndExpansion
%

%TCIMACRO{\TeXButton{Transparency}{\setbeamercovered{transparent=20}}}%
%BeginExpansion
\setbeamercovered{transparent=20}%
%EndExpansion


Una composizione $\bar{x}$ \`{e} \textbf{ammissibile} se soddisfa \emph{tutti}
i vincoli del modello: bilancio, requisito di liquidit\`{a}, limite alle
perdite, limite di concentrazione e non negativit\`{a}.

\medskip

L'insieme delle soluzioni ammissibili si indica con $\mathcal{X}$; nella
formulazione generale
\[
\mathcal{X}=\left\{  x\in\R^{n}:Ax\leq b,\;x\geq0\right\}  .
\]


La funzione obiettivo viene massimizzata \textbf{non} su tutto $\R^{n}$, ma su
$\mathcal{X}$:
\[
z^{*}=\max_{x\in\mathcal{X}}c^{\prime}x.
\]


\medskip

\begin{stepitemize}
\item Il rispetto di un \emph{singolo} requisito non rende ammissibile una
decisione: una composizione pu\`{o} avere liquidit\`{a} sufficiente e produrre
nello stesso tempo una perdita eccessiva sotto stress.

\item Una decisione con risultato economico molto elevato ma esterna a
$\mathcal{X}$ non \`{e} un'alternativa migliore: \`{e} una decisione che il
modello considera \textbf{non realizzabile}.
\end{stepitemize}

%

%TCIMACRO{\TeXButton{Transition: Box Out}{\transboxout}}%
%BeginExpansion
\transboxout
%EndExpansion%
%TCIMACRO{\TeXButton{EndFrame}{\end{frame}}}%
%BeginExpansion
\end{frame}%
%EndExpansion


%*****************************


\subsection{Ammissibilit\`{a} e confronto tra decisioni}%

%TCIMACRO{\TeXButton{BeginFrame}{\begin{frame}}}%
%BeginExpansion
\begin{frame}%
%EndExpansion
%

%TCIMACRO{\QTR{frametitle}{Il risultato pi\`{u} alto non basta}}%
%BeginExpansion
\frametitle{Il risultato pi\`{u} alto non basta}%
%EndExpansion
%

%TCIMACRO{\TeXButton{Transparency}{\setbeamercovered{transparent=20}}}%
%BeginExpansion
\setbeamercovered{transparent=20}%
%EndExpansion


Con $A_{0}=100$, $Q_{\min}=55$, $K=7$, $\alpha A_{0}=30$:

\medskip

\begin{center}
{\small
\begin{tabular}
[c]{lcccc}\hline
& \textbf{Risultato} & \textbf{Liquidit\`{a}} & \textbf{Perdita} &
\textbf{Titoli lunghi}\\
&  & $\geq55$ & $\leq7$ & $\leq30$\\\hline
$x^{(A)}=(5,10,45,40)^{\prime}$ & $4.10$ & $44.50$ & $10.15$ & $45$\\
$x^{(B)}=(10,30,20,40)^{\prime}$ & $3.65$ & $55.00$ & $6.80$ & $20$\\\hline
\end{tabular}
}
\end{center}

\medskip

\begin{stepitemize}
\item $x^{(A)}$ rende di pi\`{u} e viola \textbf{tre} requisiti su tre. Non
\`{e} un'alternativa migliore: \`{e} fuori da $\mathcal{X}$.

\item La funzione obiettivo ordina le decisioni, ma ordina soltanto quelle
ammissibili. L'\textbf{ammissibilit\`{a} precede il confronto dei valori}.

\item Le due composizioni non si distinguono per\`{o} solo in ammissibile/non
ammissibile: $x^{(B)}$ soddisfa la liquidit\`{a} esattamente al limite, la
perdita con $0.20$ di margine, la concentrazione con $10$.
\end{stepitemize}

\medskip

\begin{center}
{\small Serve una misura di \emph{quanto} una decisione dista dai propri
limiti.}
\end{center}

%

%TCIMACRO{\TeXButton{Transition: Box Out}{\transboxout}}%
%BeginExpansion
\transboxout
%EndExpansion%
%TCIMACRO{\TeXButton{EndFrame}{\end{frame}}}%
%BeginExpansion
\end{frame}%
%EndExpansion


%*****************************


\subsection{Margini, vincoli attivi, diagnosi}%

%TCIMACRO{\TeXButton{BeginFrame}{\begin{frame}}}%
%BeginExpansion
\begin{frame}%
%EndExpansion
%

%TCIMACRO{\QTR{frametitle}{Margine di un vincolo e vincoli attivi}}%
%BeginExpansion
\frametitle{Margine di un vincolo e vincoli attivi}%
%EndExpansion
%

%TCIMACRO{\TeXButton{Transparency}{\setbeamercovered{transparent=20}}}%
%BeginExpansion
\setbeamercovered{transparent=20}%
%EndExpansion


Il \textbf{margine del vincolo} $i$ misura quanto la decisione dista dal
limite del vincolo $i$. L'espressione dipende dal verso:
\[
a_{i}^{\prime}x\leq b_{i}: \quad b_{i}-a_{i}^{\prime}x; \qquad a_{i}^{\prime
}x\geq b_{i}: \quad a_{i}^{\prime}x-b_{i}.
\]


Nel primo caso $s_{i}$ conta quanto resta prima di superare il limite, nel
secondo di quanto lo si supera. In entrambi,
\[
x\in\mathcal{X} \quad\Longleftrightarrow\quad s_{i}\geq0 \;\text{ per ogni
}i.
\]


Il vincolo \`{e} \textbf{attivo} quando il suo margine \`{e} nullo, cio\`{e}
quando \`{e} soddisfatto come uguaglianza $a_{i}^{\prime}x=b_{i}$: la
decisione si trova sulla sua frontiera.

\medskip

\begin{stepitemize}
%\item I vincoli attivi sono quelli che stanno \textbf{determinando} la
%decisione. Un vincolo con margine positivo, in quel punto, non morde.


\item I vincoli attivi sono quelli che \textbf{possono contribuire a
determinare} la decisione. Un vincolo con margine positivo, in quel punto, non morde.

\item Un vincolo di uguaglianza ha margine nullo per costruzione: \`{e} attivo
in ogni soluzione ammissibile.
\end{stepitemize}

%

%TCIMACRO{\TeXButton{Transition: Box Out}{\transboxout}}%
%BeginExpansion
\transboxout
%EndExpansion%
%TCIMACRO{\TeXButton{EndFrame}{\end{frame}}}%
%BeginExpansion
\end{frame}%
%EndExpansion


%*********************


\subsection{I tre esiti della risoluzione}%

%TCIMACRO{\TeXButton{BeginFrame}{\begin{frame}}}%
%BeginExpansion
\begin{frame}%
%EndExpansion
%

%TCIMACRO{\QTR{frametitle}{I tre possibili esiti della risoluzione}}%
%BeginExpansion
\frametitle{I tre possibili esiti della risoluzione}%
%EndExpansion
%

%TCIMACRO{\TeXButton{Transparency}{\setbeamercovered{transparent=20}}}%
%BeginExpansion
\setbeamercovered{transparent=20}%
%EndExpansion


\begin{stepenumerate}
\item[1.] Una \textbf{soluzione ottima} e il corrispondente valore obiettivo.

\item[2.] L'\textbf{inammissibilit\`{a}} del modello, cio\`{e} $\mathcal{X}%
=\varnothing$: nessuna decisione soddisfa simultaneamente tutti i vincoli.

\item[3.] L'\textbf{illimitatezza}: esistono decisioni ammissibili che
producono valori arbitrariamente elevati. In un problema finanziario segnala
normalmente che manca un vincolo rilevante.
\end{stepenumerate}

\medskip

Una regione vuota pu\`{o} segnalare requisiti economicamente incompatibili,
\emph{oppure} dati e vincoli formulati in modo incoerente.

\medskip

\begin{center}
{\small \emph{Di fronte a un modello inammissibile, la prima operazione non
\`{e} modificare i vincoli fino a ottenere una soluzione.}}
\end{center}

%

%TCIMACRO{\TeXButton{Transition: Box Out}{\transboxout}}%
%BeginExpansion
\transboxout
%EndExpansion%
%TCIMACRO{\TeXButton{EndFrame}{\end{frame}}}%
%BeginExpansion
\end{frame}%
%EndExpansion


%*****************************


\section{Esercizio in aula: verifica di ammissibilit\`{a}}

%*****************************


\subsection{Enunciato}%

%TCIMACRO{\TeXButton{BeginFrame}{\begin{frame}}}%
%BeginExpansion
\begin{frame}%
%EndExpansion
%

%TCIMACRO{\QTR{frametitle}{Esercizio: questa composizione \`{e} ammissibile?}}%
%BeginExpansion
\frametitle{Esercizio: questa composizione \`{e} ammissibile?}%
%EndExpansion
%

%TCIMACRO{\TeXButton{Transparency}{\setbeamercovered{transparent=20}}}%
%BeginExpansion
\setbeamercovered{transparent=20}%
%EndExpansion


Modello bancario con $A_{0}=100$, $Q_{\min}=55$, $K=7$, $\alpha A_{0}=30$. Si
consideri la composizione
\[
x=(8,\;32,\;18,\;42)^{\prime}.
\]


\medskip

\begin{stepenumerate}


\item[1.] Verificare il vincolo di bilancio.

\item[2.] Calcolare la liquidit\`{a} disponibile e la perdita nello scenario
di stress.

\item[3.] Stabilire se la composizione \`{e} ammissibile e, in caso negativo,
quale requisito \`{e} violato.

\item[4.] Per ciascun requisito rispettato, calcolare il margine e dire se il
vincolo \`{e} attivo.
\end{stepenumerate}

\medskip

\begin{center}
$q=(1.00,\;0.90,\;0.50,\;0.20)^{\prime}$ \qquad$\ell
=(0,\;0.02,\;0.15,\;0.08)^{\prime}$
\end{center}

%

%TCIMACRO{\TeXButton{Transition: Box Out}{\transboxout}}%
%BeginExpansion
\transboxout
%EndExpansion%
%TCIMACRO{\TeXButton{EndFrame}{\end{frame}}}%
%BeginExpansion
\end{frame}%
%EndExpansion


%*****************************


\subsection{Il calcolo}%

%TCIMACRO{\TeXButton{BeginFrame}{\begin{frame}}}%
%BeginExpansion
\begin{frame}%
%EndExpansion
%

%TCIMACRO{\QTR{frametitle}{Le tre grandezze prodotte dalla composizione}}%
%BeginExpansion
\frametitle{Le tre grandezze prodotte dalla composizione}%
%EndExpansion
%

%TCIMACRO{\TeXButton{Transparency}{\setbeamercovered{transparent=20}}}%
%BeginExpansion
\setbeamercovered{transparent=20}%
%EndExpansion


Bilancio:
\[
8+32+18+42=100=A_{0}.
\]


Liquidit\`{a} disponibile:
\[
1.00(8)+0.90(32)+0.50(18)+0.20(42)=8+28.8+9+8.4=54.2.
\]


Perdita nello scenario di stress:
\[
0.02(32)+0.15(18)+0.08(42)=0.64+2.70+3.36=6.70.
\]


\medskip

\begin{stepitemize}


\item Il bilancio \`{e} un'uguaglianza: \`{e} soddisfatto esattamente da ogni
composizione che impieghi tutte le risorse.

\item La cassa non contribuisce alla perdita, perch\'{e} $\ell_{1}=0$: compare
nella liquidit\`{a} con coefficiente unitario e nella perdita non compare affatto.
\end{stepitemize}

%

%TCIMACRO{\TeXButton{Transition: Box Out}{\transboxout}}%
%BeginExpansion
\transboxout
%EndExpansion%
%TCIMACRO{\TeXButton{EndFrame}{\end{frame}}}%
%BeginExpansion
\end{frame}%
%EndExpansion


%*****************************


\subsection{La diagnosi}%

%TCIMACRO{\TeXButton{BeginFrame}{\begin{frame}}}%
%BeginExpansion
\begin{frame}%
%EndExpansion
%

%TCIMACRO{\QTR{frametitle}{Valore, limite, margine}}%
%BeginExpansion
\frametitle{Valore, limite, margine}%
%EndExpansion
%

%TCIMACRO{\TeXButton{Transparency}{\setbeamercovered{transparent=20}}}%
%BeginExpansion
\setbeamercovered{transparent=20}%
%EndExpansion


\begin{center}
{\small
\begin{tabular}
[c]{lcccc}\hline
\textbf{Vincolo} & \textbf{Valore} & \textbf{Limite} & \textbf{Margine} &
\textbf{Stato}\\\hline
Bilancio & $100$ & $=100$ & $0$ & attivo\\
Liquidit\`{a} & $54.2$ & $\geq55$ & $-0.8$ & violato\\
Perdita sotto stress & $6.70$ & $\leq7$ & $0.30$ & non attivo\\
Concentrazione ($x_{3}$) & $18$ & $\leq30$ & $12$ & non attivo\\\hline
\end{tabular}
}
\end{center}

\medskip

\begin{stepitemize}
\item Un solo requisito violato basta: $x\notin\mathcal{X}$. Il risultato
economico della composizione non viene nemmeno calcolato.

\item Tra i vincoli rispettati \`e attivo soltanto il bilancio, come ogni
uguaglianza; perdita e concentrazione conservano margine.

\item Lo scostamento $-0.8$ misura quanta liquidit\`{a} manca. \`{E} la
quantit\`{a} che una correzione dovrebbe colmare.
\end{stepitemize}

%

%TCIMACRO{\TeXButton{Transition: Box Out}{\transboxout}}%
%BeginExpansion
\transboxout
%EndExpansion%
%TCIMACRO{\TeXButton{EndFrame}{\end{frame}}}%
%BeginExpansion
\end{frame}%
%EndExpansion


%*****************************


\section{La geometria del problema}

%*****************************


\subsection{Un modello bancario ridotto}%

%TCIMACRO{\TeXButton{BeginFrame}{\begin{frame}}}%
%BeginExpansion
\begin{frame}%
%EndExpansion
%

%TCIMACRO{\QTR{frametitle}{Un problema bancario ridotto, rappresentabile
%graficamente}}%
%BeginExpansion
\frametitle{Un problema bancario ridotto, rappresentabile graficamente}%
%EndExpansion
%

%TCIMACRO{\TeXButton{Transparency}{\setbeamercovered{transparent=20}}}%
%BeginExpansion
\setbeamercovered{transparent=20}%
%EndExpansion


Il modello a quattro classi non pu\`{o} essere rappresentato su un piano.
Consideriamo una versione ridotta a \textbf{due sole grandezze}: le
attivit\`{a} liquide e i titoli a lunga scadenza.

\medskip

\begin{center}
{\small
\begin{tabular}
[c]{ll}\hline
$x_{1}$ & attivit\`{a} liquide\\
$x_{2}$ & titoli a lunga scadenza\\\hline
\end{tabular}
}

{\small \emph{Avvertenza: in questa rappresentazione le due variabili sono
ridefinite. }}
\end{center}

\medskip

\begin{stepitemize}
\item Le risorse destinate alle due classi non possono superare $A_{0}$: la
parte restante \`{e} investita nelle altre classi, non rappresentate.

\item Restano un requisito minimo sulle liquide e un limite massimo sulle
lunghe, oltre alla non negativit\`{a}.
\end{stepitemize}

%

\[
x_{1}+x_{2}\leq10, \qquad x_{1}\geq3, \qquad x_{2}\leq4, \qquad x_{1}%
,x_{2}\geq0.
\]
%

%TCIMACRO{\TeXButton{Transition: Box Out}{\transboxout}}%
%BeginExpansion
\transboxout
%EndExpansion%
%TCIMACRO{\TeXButton{EndFrame}{\end{frame}}}%
%BeginExpansion
\end{frame}%
%EndExpansion


%*****************************


\subsection{Regione ammissibile e vertici}%

%TCIMACRO{\TeXButton{BeginFrame}{\begin{frame}}}%
%BeginExpansion
\begin{frame}%
%EndExpansion
%

%TCIMACRO{\QTR{frametitle}{La regione ammissibile e i suoi vertici}}%
%BeginExpansion
\frametitle{La regione ammissibile e i suoi vertici}%
%EndExpansion
%

%TCIMACRO{\TeXButton{Transparency}{\setbeamercovered{transparent=20}}}%
%BeginExpansion
\setbeamercovered{transparent=20}%
%EndExpansion


\begin{center}%
%TCIMACRO{\FRAME{itbpF}{2.9365in}{1.9561in}{0in}{}{}%
%{cap11_geometria_programmazione_lineare.png}%
%{\special{ language "Scientific Word";  type "GRAPHIC";
%maintain-aspect-ratio TRUE;  display "USEDEF";  valid_file "F";
%width 2.9365in;  height 1.9561in;  depth 0in;  original-width 8.4567in;
%original-height 5.6333in;  cropleft "0";  croptop "1";  cropright "1";
%cropbottom "0";
%filename '../graphics/Cap11_geometria_programmazione_lineare.png';file-properties "XNPEU";}%
%} }%
%BeginExpansion
{\includegraphics[
natheight=5.633300in,
natwidth=8.456700in,
height=1.9561in,
width=2.9365in
]%
{../graphics/Cap11_geometria_programmazione_lineare.png}%
}
%EndExpansion


{\small
\begin{tabular}
[c]{ccl}\hline
\textbf{Vertice} & \textbf{Coordinate} & \textbf{Vincoli attivi}\\\hline
$A$ & $(3,0)$ & liquide al minimo, nessuna lunga\\
$B$ & $(10,0)$ & risorse esaurite, nessuna lunga\\
$C$ & $(6,4)$ & risorse esaurite, lunghe al massimo\\
$D$ & $(3,4)$ & liquide al minimo, lunghe al massimo\\\hline
\end{tabular}
}
\end{center}

%

%TCIMACRO{\TeXButton{Transition: Box Out}{\transboxout}}%
%BeginExpansion
\transboxout
%EndExpansion%
%TCIMACRO{\TeXButton{EndFrame}{\end{frame}}}%
%BeginExpansion
\end{frame}%
%EndExpansion


%*****************************


\subsection{Rette di livello e punti estremi}%

%TCIMACRO{\TeXButton{BeginFrame}{\begin{frame}}}%
%BeginExpansion
\begin{frame}%
%EndExpansion
%

%TCIMACRO{\QTR{frametitle}{Come la funzione obiettivo sceglie un vertice}}%
%BeginExpansion
\frametitle{Come la funzione obiettivo sceglie un vertice}%
%EndExpansion
%

%TCIMACRO{\TeXButton{Transparency}{\setbeamercovered{transparent=20}}}%
%BeginExpansion
\setbeamercovered{transparent=20}%
%EndExpansion


Stessa regione ammissibile, due criteri diversi: $2x_{1}+3x_{2}$ e
$3x_{1}+2x_{2}$.

\begin{center}%
%TCIMACRO{\FRAME{itbpF}{2.8418in}{2.0845in}{0in}{}{}%
%{cap11_es02_geometria_ottimo.png}{\special{ language "Scientific Word";
%type "GRAPHIC";  maintain-aspect-ratio TRUE;  display "USEDEF";
%valid_file "F";  width 2.8418in;  height 2.0845in;  depth 0in;
%original-width 7.1300in;  original-height 5.2300in;  cropleft "0";
%croptop "1";  cropright "1";  cropbottom "0";
%filename '../graphics/Cap11_Es02_geometria_ottimo.png';file-properties "XNPEU";}%
%} }%
%BeginExpansion
{\includegraphics[
natheight=5.230000in,
natwidth=7.130000in,
height=2.0845in,
width=2.8418in
]%
{../graphics/Cap11_Es02_geometria_ottimo.png}%
}
%EndExpansion

\end{center}

\begin{stepitemize}
\item Fissato un livello $k$, l'equazione $c_{1}x_{1}+c_{2}x_{2}=k$ descrive
una \textbf{retta di livello}: tutti i suoi punti danno lo stesso risultato.
Al crescere di $k$ la retta trasla parallelamente.
\end{stepitemize}

%

%TCIMACRO{\TeXButton{Transition: Box Out}{\transboxout}}%
%BeginExpansion
\transboxout
%EndExpansion%
%TCIMACRO{\TeXButton{EndFrame}{\end{frame}}}%
%BeginExpansion
\end{frame}%
%EndExpansion


%*****************************


\subsection{Vincoli attivi e vertici}%

%TCIMACRO{\TeXButton{BeginFrame}{\begin{frame}}}%
%BeginExpansion
\begin{frame}%
%EndExpansion
%

%TCIMACRO{\QTR{frametitle}{Perch\'{e} l'ottimo si trova dove i vincoli
%stringono}}%
%BeginExpansion
\frametitle{Perch\'{e} l'ottimo si trova dove i vincoli stringono}%
%EndExpansion
%

%TCIMACRO{\TeXButton{Transparency}{\setbeamercovered{transparent=20}}}%
%BeginExpansion
\setbeamercovered{transparent=20}%
%EndExpansion


%Un \textbf{vertice} \`{e} il punto in cui si incontrano tante frontiere quante sono le variabili. In due dimensioni servono due vincoli attivi; con $n$ variabili ne servono $n$.


Un \textbf{vertice} \`{e} un punto della regione ammissibile nel quale sono
attivi vincoli sufficienti a determinare localmente la soluzione. Con $n$
variabili, servono $n$ vincoli attivi \textbf{linearmente indipendenti} per
determinare il vertice; in casi degeneri, possono essere attivi anche pi\`{u}
di $n$ vincoli.

\begin{stepitemize}
\item L'ottimo \`{e} l'\textbf{ultimo punto di contatto} con la regione. I
vincoli non cambiano, cambia l'inclinazione delle rette: l'ottimo si sposta da
$C$ a $B$.

\item Se le rette fossero parallele a un lato, tutti i punti di quel lato
sarebbero ottimi. In ogni caso \textbf{almeno un vertice \`{e} ottimo}: per
cercare l'ottimo basta guardare i vertici.
\end{stepitemize}

\medskip

\begin{center}
{\small Cercare l'ottimo significa cercare la \textbf{combinazione di vincoli}
che lo determina. Il valore di $z^{*}$ dice quanto si ottiene; i vincoli
attivi dicono \emph{perch\'{e}} non si ottiene di pi\`{u}.}
\end{center}

%

%TCIMACRO{\TeXButton{Transition: Box Out}{\transboxout}}%
%BeginExpansion
\transboxout
%EndExpansion%
%TCIMACRO{\TeXButton{EndFrame}{\end{frame}}}%
%BeginExpansion
\end{frame}%
%EndExpansion


%*****************************


\subsection{Che cosa significa soluzione ottima}%

%TCIMACRO{\TeXButton{BeginFrame}{\begin{frame}}}%
%BeginExpansion
\begin{frame}%
%EndExpansion
%

%TCIMACRO{\QTR{frametitle}{Che cosa significa soluzione ottima}}%
%BeginExpansion
\frametitle{Che cosa significa soluzione ottima}%
%EndExpansion
%

%TCIMACRO{\TeXButton{Transparency}{\setbeamercovered{transparent=20}}}%
%BeginExpansion
\setbeamercovered{transparent=20}%
%EndExpansion


Una soluzione $x^{*}\in\mathcal{X}$ \`{e} \textbf{ottima} se
\[
c^{\prime}x^{*}\geq c^{\prime}x \qquad\text{per ogni }x\in\mathcal{X},
\]
e il valore $z^{*}=c^{\prime}x^{*}$ \`{e} il \textbf{valore ottimo}. La
definizione richiede due condizioni: ammissibilit\`{a}, e nessuna alternativa
ammissibile migliore.

\medskip

\begin{stepitemize}
\item Quanto visto nel piano vale in generale: se $\mathcal{X}$ non \`e vuoto
e il massimo \`e finito, almeno una soluzione ottima \`e un \textbf{punto
estremo}.

\item Tra i vincoli attivi che individuano un vertice possono figurare le
condizioni di non negativit\`a, che sono vincoli a tutti gli effetti.
\end{stepitemize}

\medskip

\begin{center}
{\small L'ottimo va riferito al modello adottato: variabili incluse,
coefficienti, vincoli, funzione obiettivo. Non \`{e} la migliore decisione
possibile in senso assoluto.}
\end{center}

%

%TCIMACRO{\TeXButton{Transition: Box Out}{\transboxout}}%
%BeginExpansion
\transboxout
%EndExpansion%
%TCIMACRO{\TeXButton{EndFrame}{\end{frame}}}%
%BeginExpansion
\end{frame}%
%EndExpansion


%*****************************


\section{La soluzione ottima del modello SVB}

\subsection{La composizione ottima}%

%TCIMACRO{\TeXButton{BeginFrame}{\begin{frame}}}%
%BeginExpansion
\begin{frame}%
%EndExpansion
%

%TCIMACRO{\QTR{frametitle}{Torniamo al modello completo SVB: soluzione
%ottima}}%
%BeginExpansion
\frametitle{Torniamo al modello completo SVB: soluzione ottima}%
%EndExpansion
%

%TCIMACRO{\TeXButton{Transparency}{\setbeamercovered{transparent=20}}}%
%BeginExpansion
\setbeamercovered{transparent=20}%
%EndExpansion


La soluzione ottima del caso \`{e} determinata da quattro vincoli attivi
linearmente indipendenti. Risolvendo con $A_{0}=100$, $Q_{\min}=55$, $K=7$ e
$x_{3}\leq30$, si ottiene $x^{*}\approx(0;\;41.045;\;20.896;\;38.060)^{\prime
}$, $z^{*}\approx3.7649$.

\begin{center}%
%TCIMACRO{\FRAME{itbpF}{2.6050in}{1.4409in}{0in}{}{}%
%{cap11_soluzione_ottima_allocazione.png}%
%{\special{ language "Scientific Word";  type "GRAPHIC";
%maintain-aspect-ratio TRUE;  display "USEDEF";  valid_file "F";
%width 2.6050in;  height 1.4409in;  depth 0in;  original-width 11.5955in;
%original-height 6.4136in;  cropleft "0";  croptop "1";  cropright "1";
%cropbottom "0";
%filename '../graphics/Cap11_soluzione_ottima_allocazione.png';file-properties "XNPEU";}%
%} }%
%BeginExpansion
{\includegraphics[
natheight=6.413600in,
natwidth=11.595500in,
height=1.4409in,
width=2.6050in
]%
{../graphics/Cap11_soluzione_ottima_allocazione.png}%
}
%EndExpansion



\end{center}

\vspace{-0.1cm}

\begin{stepitemize}
\item La soluzione \`{e} un vettore, non un numero: $z^{\ast}$ dice
\textbf{quanto} si ottiene, $x^{\ast}$ dice \textbf{come}. Due modelli con lo
stesso valore ottimo possono richiedere decisioni molto diverse.

\item $z^{*}$ ha l'unit\`{a} dei coefficienti $c_{j}$: \`{e} il margine
complessivo, che con $A_{0}=100$ si legge anche come percentuale dell'attivo.
\end{stepitemize}

%

%TCIMACRO{\TeXButton{Transition: Box Out}{\transboxout}}%
%BeginExpansion
\transboxout
%EndExpansion%
%TCIMACRO{\TeXButton{EndFrame}{\end{frame}}}%
%BeginExpansion
\end{frame}%
%EndExpansion


%*****************************


\subsection{Che cosa rivelano i vincoli attivi}%

%TCIMACRO{\TeXButton{BeginFrame}{\begin{frame}}}%
%BeginExpansion
\begin{frame}%
%EndExpansion


%\QTR{frametitle}{I vincoli attivi identificano la soluzione}
%

%TCIMACRO{\QTR{frametitle}{Che cosa rivelano i vincoli attivi}}%
%BeginExpansion
\frametitle{Che cosa rivelano i vincoli attivi}%
%EndExpansion
%

%TCIMACRO{\TeXButton{Transparency}{\setbeamercovered{transparent=20}}}%
%BeginExpansion
\setbeamercovered{transparent=20}%
%EndExpansion


\begin{center}
{\small
\begin{tabular}
[c]{lcccc}\hline
\textbf{Vincolo} & \textbf{Valore} & \textbf{Limite} & \textbf{Margine} &
\textbf{Stato}\\\hline
Bilancio & $100$ & $=100$ & $0$ & attivo\\
Liquidit\`{a} & $55.00$ & $\geq55$ & $0$ & attivo\\
Perdita sotto stress & $7.00$ & $\leq7$ & $0$ & attivo\\
Non negativit\`{a} di $x_{1}$ & $0$ & $\geq0$ & $0$ & attivo\\
Concentrazione ($x_{3}$) & $20.896$ & $\leq30$ & $9.104$ & non attivo\\\hline
\end{tabular}
}
\end{center}

\medskip

Quattro variabili, quattro vincoli attivi: la soluzione \`{e} il vertice ottimo.

\medskip

\begin{stepitemize}
\item Un vincolo con margine \textbf{positivo} non condiziona localmente la
soluzione: un piccolo allentamento lascia invariati $x^{*}$ e $z^{*}$.

%\item Un vincolo con margine \textbf{positivo} non condiziona la soluzione:
%allentarlo lascerebbe $x^{*}$ invariato.


\item Solo un vincolo \textbf{attivo} pu\`{o} avere valore marginale non
nullo. La decisione \`{e} condizionata dalla combinazione tra fabbisogno di
liquidit\`{a} e capacit\`{a} di assorbire perdite.

\item Un vincolo attivo pu\`{o} rivelare un requisito \textbf{mancante}:
$x_{1}^{*}=0$ non significa cassa inutile, ma che nessun vincolo impone di
detenerne. Se la banca ritiene necessaria una dotazione minima, deve scriverla
come vincolo: $x_{1}\geq C_{\min}$.
\end{stepitemize}

%

%TCIMACRO{\TeXButton{Transition: Box Out}{\transboxout}}%
%BeginExpansion
\transboxout
%EndExpansion%
%TCIMACRO{\TeXButton{EndFrame}{\end{frame}}}%
%BeginExpansion
\end{frame}%
%EndExpansion


%**************************


\section{Sensitivit\`a e prezzi ombra}

\subsection{La funzione del valore ottimo}%

%TCIMACRO{\TeXButton{BeginFrame}{\begin{frame}}}%
%BeginExpansion
\begin{frame}%
%EndExpansion
%

%TCIMACRO{\QTR{frametitle}{Il valore ottimo in funzione del requisito di
%liquidit\`a}}%
%BeginExpansion
\frametitle{Il valore ottimo in funzione del requisito di liquidit\`a}%
%EndExpansion
%

%TCIMACRO{\TeXButton{Transparency}{\setbeamercovered{transparent=20}}}%
%BeginExpansion
\setbeamercovered{transparent=20}%
%EndExpansion


Come ogni termine noto, il requisito di liquidit\`a \`e fissato prima della
soluzione: finora $Q_{\min}=55$. Ora lo trattiamo come un \textbf{parametro}
$Q$ e chiamiamo
\[%
\begin{array}
[c]{rl}%
\mathcal{X}(Q) & \text{le composizioni ammissibili con il vincolo }q^{\prime
}x\geq Q,\\[3pt]%
z^{*}(Q) & \text{il miglior risultato ottenibile in }\mathcal{X}(Q).
\end{array}
\]


\begin{stepitemize}
\item \textbf{Un requisito pi\`u severo restringe le scelte.} Se $Q_{2}>Q_{1}%
$, ogni composizione con $q^{\prime}x\geq Q_{2}$ ha anche $q^{\prime}x\geq
Q_{1}$, quindi $\mathcal{X}(Q_{2})\subseteq\mathcal{X}(Q_{1})$.

\item \textbf{Ridurre lo spazio delle soluzioni ammissibili non pu\`{o}
migliorare il risultato.} La composizione ottima con $Q_{2}$ vale altrettanto
con $Q_{1}$, quindi $z^{\ast}(Q_{2})\leq z^{\ast}(Q_{1})$.

\item \textbf{La differenza \`e il costo della prudenza.} $z^{*}(Q_{1}%
)-z^{*}(Q_{2})$ \`e il rendimento a cui si rinuncia per avere pi\`u
liquidit\`a: un \textbf{costo opportunit\`a}, misurato \emph{all'interno del
modello}. Non dice che la liquidit\`a sia un male: dice quanto costa.
\end{stepitemize}

%

%TCIMACRO{\TeXButton{Transition: Box Out}{\transboxout}}%
%BeginExpansion
\transboxout
%EndExpansion%
%TCIMACRO{\TeXButton{EndFrame}{\end{frame}}}%
%BeginExpansion
\end{frame}%
%EndExpansion


%*****************************


\subsection{Analisi numerica e costo marginale}%

%TCIMACRO{\TeXButton{BeginFrame}{\begin{frame}}}%
%BeginExpansion
\begin{frame}%
%EndExpansion
%

%TCIMACRO{\QTR{frametitle}{Il costo marginale della liquidit\`{a}}}%
%BeginExpansion
\frametitle{Il costo marginale della liquidit\`{a}}%
%EndExpansion
%

%TCIMACRO{\TeXButton{Transparency}{\setbeamercovered{transparent=20}}}%
%BeginExpansion
\setbeamercovered{transparent=20}%
%EndExpansion


\begin{center}
{\small
\begin{tabular}
[c]{ccl}\hline
$Q$ & $z^{*}(Q)$ & \textbf{Composizione ottima}\\\hline
$55$ & $3.7649$ & $(0;\;41.045;\;20.896;\;38.060)$\\
$65$ & $3.4142$ & $(0;\;51.493;\;29.851;\;18.657)$\\
$75$ & $3.0571$ & $(0;\;65.714;\;30.000;\;4.286)$\\
$80$ & $2.8750$ & $(0;\;75.000;\;25.000;\;0)$\\\hline
\end{tabular}
}

\bigskip%

%TCIMACRO{\FRAME{itbpF}{2.6368in}{1.5865in}{0in}{}{}%
%{cap11_valore_ottimo_liquidita.png}{\special{ language "Scientific Word";
%type "GRAPHIC";  maintain-aspect-ratio TRUE;  display "USEDEF";
%valid_file "F";  width 2.6368in;  height 1.5865in;  depth 0in;
%original-width 6.7933in;  original-height 4.06in;  cropleft "0";
%croptop "1";  cropright "1";  cropbottom "0";
%filename '../graphics/Cap11_valore_ottimo_liquidita.png';file-properties "XNPEU";}%
%} }%
%BeginExpansion
{\includegraphics[
natheight=4.060000in,
natwidth=6.793300in,
height=1.5865in,
width=2.6368in
]%
{../graphics/Cap11_valore_ottimo_liquidita.png}%
}
%EndExpansion



\end{center}

$z^{\ast}(Q)$ \`{e} \textbf{funzione lineare a tratti}. La pendenza cambia
quando cambia l'insieme dei vincoli attivi.%

%TCIMACRO{\TeXButton{Transition: Box Out}{\transboxout}}%
%BeginExpansion
\transboxout
%EndExpansion%
%TCIMACRO{\TeXButton{EndFrame}{\end{frame}}}%
%BeginExpansion
\end{frame}%
%EndExpansion


%*****************************


\subsection{Una valutazione locale}%

%TCIMACRO{\TeXButton{BeginFrame}{\begin{frame}}}%
%BeginExpansion
\begin{frame}%
%EndExpansion
%

%TCIMACRO{\QTR{frametitle}{Come si legge un costo marginale}}%
%BeginExpansion
\frametitle{Come si legge un costo marginale}%
%EndExpansion
%

%TCIMACRO{\TeXButton{Transparency}{\setbeamercovered{transparent=20}}}%
%BeginExpansion
\setbeamercovered{transparent=20}%
%EndExpansion


\begin{stepitemize}
\item \textbf{\`{E} un costo opportunit\`{a}, non una spesa.} Nessun esborso
contabile: \`{e} il risultato massimo al quale il modello rinuncia per
rispettare un requisito leggermente pi\`{u} severo.

\item \textbf{Vale localmente.} La derivata di $z^{*}$ rispetto al termine
noto \`{e} costante finch\'{e} l'insieme dei vincoli attivi non cambia. Fuori
da quell'intervallo va ricalcolata: un costo marginale senza il proprio
intervallo di validit\`{a} \`{e} un numero incompleto.

\item \textbf{Si legge nella ricomposizione.} Il costo non nasce dal nulla:
irrigidire un requisito obbliga a sostituire alcune attivit\`{a} con altre
meno redditizie. Qui le risorse migrano verso i titoli a breve e i prestiti si
riducono fino ad azzerarsi.

\item \textbf{Non \`{e} una risposta normativa.} Dice quanto costa la
prudenza, non quanta prudenza serva: quella scelta richiede valutazioni sul
rischio dei deflussi e sulla stabilit\`{a} della raccolta, che il modello non contiene.
\end{stepitemize}

%

%TCIMACRO{\TeXButton{Transition: Box Out}{\transboxout}}%
%BeginExpansion
\transboxout
%EndExpansion%
%TCIMACRO{\TeXButton{EndFrame}{\end{frame}}}%
%BeginExpansion
\end{frame}%
%EndExpansion


%*****************************


\subsection{Valore marginale dei vincoli e prezzi ombra}%

%TCIMACRO{\TeXButton{BeginFrame}{\begin{frame}}}%
%BeginExpansion
\begin{frame}%
%EndExpansion
%

%TCIMACRO{\QTR{frametitle}{Il prezzo ombra di un vincolo}}%
%BeginExpansion
\frametitle{Il prezzo ombra di un vincolo}%
%EndExpansion
%

%TCIMACRO{\TeXButton{Transparency}{\setbeamercovered{transparent=20}}}%
%BeginExpansion
\setbeamercovered{transparent=20}%
%EndExpansion


Per la liquidit\`{a} di SVB abbiamo visto che: requisito pi\`{u} severo
$\Rightarrow$ insieme ammissibile pi\`{u} piccolo $\Rightarrow$ valore ottimo
inferiore. Questo ragionamento si estende ad \textbf{ogni} vincolo attivo ed
introduce il concetto di \textbf{prezzo ombra}.

\begin{definition}
[Definizione]Il \textbf{prezzo ombra} $\lambda_{i}$ di un vincolo $i$ \`{e} la
derivata del valore ottimo rispetto al termine noto $b_{i}$:
\[
\lambda_{i}=\frac{\partial z^{\ast}}{\partial b_{i}}.
\]

\end{definition}

\begin{itemize}
\item Finch\'{e} i vincoli attivi non cambiano, $z^{\ast}$ varia
\textbf{linearmente} con ciascun termine noto e quindi $\lambda_{i}$ =
\textbf{valore costante} (come abbiamo visto per $Q$ nel caso SVB).

\item Un vincolo con margine positivo non condiziona la soluzione: una piccola
variazione del suo termine noto lascia invariati $x^{\ast}$ e $z^{\ast}$.
\textbf{Solo i vincoli attivi hanno un valore marginale diverso da zero.}
\end{itemize}

%

%TCIMACRO{\TeXButton{Transition: Box Out}{\transboxout}}%
%BeginExpansion
\transboxout
%EndExpansion%
%TCIMACRO{\TeXButton{EndFrame}{\end{frame}}}%
%BeginExpansion
\end{frame}%
%EndExpansion


%*****************************
%

%TCIMACRO{\TeXButton{BeginFrame}{\begin{frame}}}%
%BeginExpansion
\begin{frame}%
%EndExpansion
%

%TCIMACRO{\QTR{frametitle}{Ogni vincolo attivo ha un valore marginale}}%
%BeginExpansion
\frametitle{Ogni vincolo attivo ha un valore marginale}%
%EndExpansion
%

%TCIMACRO{\TeXButton{Transparency}{\setbeamercovered{transparent=20}}}%
%BeginExpansion
\setbeamercovered{transparent=20}%
%EndExpansion


\begin{itemize}
\item \textbf{Interpretazione dei prezzi ombra}: dipende dal tipo di problema,
di massimo (es.\ un profitto) o di minimo (es.\ un costo), e dal verso in cui
$b_{i}$ modifica $\mathcal{X}$:
\end{itemize}

\begin{center}
{\small
\begin{tabular}
[c]{lcc}\hline
& aumentare $b_{i}$ \textbf{allarga} $\mathcal{X}$ & aumentare $b_{i}$
\textbf{restringe} $\mathcal{X}$\\\hline
\textbf{max} (profitto) & $\lambda_{i}>0$ & $\lambda_{i}<0$\\
& {\footnotesize profitto dal rilassamento} & {\footnotesize perdita dalla
restrizione}\\[4pt]%
\textbf{min} (costo) & $\lambda_{i}<0$ & $\lambda_{i}>0$\\
& {\footnotesize risparmio dal rilassamento} & {\footnotesize costo della
restrizione}\\\hline
\end{tabular}
}
\end{center}

\begin{itemize}
\item \textbf{Caso SVB, massimizzare rendimento, vincolo liquidit\`{a} }$dQ$
restringe $\mathcal{X}$\textbf{:} $\partial z^{\ast}/\partial Q\approx
-0.03507$. Ogni unit\`{a} di requisito liquidit\`{a} aggiuntiva $Q$
produrrebbe una riduzione del rendimento.
\end{itemize}

\medskip

{\small Calcolare i prezzi ombra un vincolo alla volta richiede di risolvere
pi\`u volte il modello. Il \textbf{problema duale} (Lezione 12) li fornisce
tutti insieme.}%

%TCIMACRO{\TeXButton{Transition: Box Out}{\transboxout}}%
%BeginExpansion
\transboxout
%EndExpansion%
%TCIMACRO{\TeXButton{EndFrame}{\end{frame}}}%
%BeginExpansion
\end{frame}%
%EndExpansion


\section{Esercizio in aula: redditivit\`{a} e resilienza}

%*****************************


\subsection{Enunciato}%

%TCIMACRO{\TeXButton{BeginFrame}{\begin{frame}}}%
%BeginExpansion
\begin{frame}%
%EndExpansion
%

%TCIMACRO{\QTR{frametitle}{Esercizio: il costo di una politica pi\`{u}
%resiliente}}%
%BeginExpansion
\frametitle{Esercizio: il costo di una politica pi\`{u} resiliente}%
%EndExpansion
%

%TCIMACRO{\TeXButton{Transparency}{\setbeamercovered{transparent=20}}}%
%BeginExpansion
\setbeamercovered{transparent=20}%
%EndExpansion


La banca valuta una politica pi\`{u} prudente. Rispetto alla soluzione ottima
trovata, cambiano due requisiti:

\begin{center}
{\small
\begin{tabular}
[c]{lcc}\hline
\textbf{Requisito} & \textbf{Caso base} & \textbf{Politica resiliente}\\\hline
Liquidit\`{a} minima $Q_{\min}$ & $55$ & $75$\\
Cassa minima $x_{1}\geq C_{\min}$ & --- & $10$\\\hline
\end{tabular}
}
\end{center}

Bilancio ($A_{0}=100$), perdita sotto stress ($K=7$) e concentrazione
($x_{3}\leq30$) restano invariati. Prima di risolvere, formulate una previsione:

\begin{stepitemize}
\item Il valore ottimo $z^{\ast}$ aumenter\`{a} o diminuir\`{a}? Si pu\`{o}
dire senza calcoli?

\item Quali vincoli resteranno attivi e quali potrebbero acquistare margine?
Partite dai margini della soluzione ottima.

\item Che segno vi aspettate per i prezzi ombra dei due requisiti irrigiditi?
\end{stepitemize}

%

%TCIMACRO{\TeXButton{Transition: Box Out}{\transboxout}}%
%BeginExpansion
\transboxout
%EndExpansion%
%TCIMACRO{\TeXButton{EndFrame}{\end{frame}}}%
%BeginExpansion
\end{frame}%
%EndExpansion


%*****************************


\subsection{Debrief}%

%TCIMACRO{\TeXButton{BeginFrame}{\begin{frame}}}%
%BeginExpansion
\begin{frame}%
%EndExpansion
%

%TCIMACRO{\QTR{frametitle}{I vincoli si scambiano il ruolo}}%
%BeginExpansion
\frametitle{I vincoli si scambiano il ruolo}%
%EndExpansion
%

%TCIMACRO{\TeXButton{Transparency}{\setbeamercovered{transparent=20}}}%
%BeginExpansion
\setbeamercovered{transparent=20}%
%EndExpansion


\begin{center}
{\small
\begin{tabular}
[c]{lrcrrc}\hline
& \multicolumn{2}{c}{\textbf{Caso base}} &  &
\multicolumn{2}{c}{\textbf{Politica resiliente}}\\\cline{2-3}\cline{5-6}%
\textbf{Vincolo} & \textbf{margine} & $\lambda_{i}$ &  & \textbf{margine} &
$\lambda_{i}$\\\hline
Liquidit\`{a} & $0$ & \multicolumn{1}{r}{$-0.03507$} &  & $0$ &
\multicolumn{1}{r}{$-0.03571$}\\
Perdita sotto stress & $0$ & \multicolumn{1}{r}{$0.00746$} &  & $0.957$ &
\multicolumn{1}{r}{$0$}\\
Concentrazione & $9.104$ & \multicolumn{1}{r}{$0$} &  & $0$ &
\multicolumn{1}{r}{$0.00071$}\\
Cassa minima & --- & \multicolumn{1}{r}{---} &  & $0$ &
\multicolumn{1}{r}{$-0.01143$}\\\hline
Soluzione ottima $x^{\ast}$ & \multicolumn{2}{c}{$(0;\,41.0;\,20.9;\,38.1)$} &
& \multicolumn{2}{c}{$(10;\,54.3;\,30;\,5.7)$}\\
Valore ottimo $z^{\ast}$ & \multicolumn{2}{c}{$3.7649$} &  &
\multicolumn{2}{c}{$2.9429$}\\\hline
\end{tabular}
}
\end{center}

\begin{stepitemize}
\item \textbf{Il costo della resilienza} \`{e} $3.7649-2.9429=0.822$, circa il
$22\%$ del valore ottimo: le risorse passano dai prestiti ai titoli a breve e
alla cassa.

\item \textbf{Alcuni vincoli attivi si sostituiscono}: la perdita sotto stress
acquista margine e il suo prezzo ombra si annulla; concentrazione e cassa
minima diventano attive. I prezzi ombra sono \textbf{locali}: cambiano con
l'insieme dei vincoli attivi.
\end{stepitemize}

%

%TCIMACRO{\TeXButton{Transition: Box Out}{\transboxout}}%
%BeginExpansion
\transboxout
%EndExpansion%
%TCIMACRO{\TeXButton{EndFrame}{\end{frame}}}%
%BeginExpansion
\end{frame}%
%EndExpansion


%*****************************


\section{Limiti del modello e sintesi}

%*****************************


\subsection{Che cosa il modello non rappresenta}%

%TCIMACRO{\TeXButton{BeginFrame}{\begin{frame}}}%
%BeginExpansion
\begin{frame}%
%EndExpansion
%

%TCIMACRO{\QTR{frametitle}{I limiti sono conseguenze delle ipotesi}}%
%BeginExpansion
\frametitle{I limiti sono conseguenze delle ipotesi}%
%EndExpansion
%

%TCIMACRO{\TeXButton{Transparency}{\setbeamercovered{transparent=20}}}%
%BeginExpansion
\setbeamercovered{transparent=20}%
%EndExpansion


Le ipotesi del modello lo rendono risolvibile, ma ciascuna esclude un fenomeno
che nel caso SVB ha avuto un ruolo.

\begin{center}
{\small
\begin{tabular}
[c]{ll}\hline
\textbf{Ipotesi} & \textbf{Che cosa esclude}\\\hline
Proporzionalit\`{a} & le vendite forzate (\emph{fire sale}): $\ell_{j}$ non
cresce\\
& con la quantit\`{a} venduta\\[3pt]%
Additivit\`{a} & il contagio: le perdite su una classe non costringono\\
& a vendere in perdita le altre\\[3pt]%
Coefficienti noti & l'incertezza: un solo scenario di stress. Pi\`{u}
scenari\\
& sono l'oggetto della programmazione stocastica\\
& (Lezioni 14--15)\\\hline
\end{tabular}
}
\end{center}

\begin{stepitemize}
\item Anche la \textbf{funzione obiettivo} ha un limite: misura solo il
margine. Il costo della resilienza \`{e} nel modello ($z^{\ast}$ scende di
$0.822$); il suo beneficio, resistere a una corsa agli sportelli come quella
subita da SVB, non lo \`{e}. Servirebbe ricorrere a \textbf{metodi
multi-criterio}.
\end{stepitemize}

%

%TCIMACRO{\TeXButton{Transition: Box Out}{\transboxout}}%
%BeginExpansion
\transboxout
%EndExpansion%
%TCIMACRO{\TeXButton{EndFrame}{\end{frame}}}%
%BeginExpansion
\end{frame}%
%EndExpansion


%*****************************


\subsection{Sintesi}%

%TCIMACRO{\TeXButton{BeginFrame}{\begin{frame}}}%
%BeginExpansion
\begin{frame}%
%EndExpansion
%

%TCIMACRO{\QTR{frametitle}{Sintesi: dal vincolo al suo prezzo}}%
%BeginExpansion
\frametitle{Sintesi: dal vincolo al suo prezzo}%
%EndExpansion
%

%TCIMACRO{\TeXButton{Transparency}{\setbeamercovered{transparent=20}}}%
%BeginExpansion
\setbeamercovered{transparent=20}%
%EndExpansion


\begin{stepenumerate}
\item[1.] \textbf{Formulare.} Variabili decisionali, funzione obiettivo e
vincoli: $\max c^{\prime}x$, $Ax\leq b$, $x\geq0$.

\item[2.] \textbf{Delimitare.} $\mathcal{X}$ raccoglie le decisioni
ammissibili; il margine misura la distanza di un vincolo dal proprio limite.

\item[3.] \textbf{Individuare.} L'ottimo pu\`{o} essere cercato tra i punti
estremi, determinati dai vincoli attivi.

\item[4.] \textbf{Valutare al margine.} La sensitivit\`{a} misura come cambia
localmente il valore ottimo al variare dei requisiti.

\item[5.] \textbf{Interpretare.} Il prezzo ombra esprime il valore marginale
di un vincolo; se il vincolo ha margine positivo, il prezzo \`{e} localmente nullo.
\end{stepenumerate}

\begin{center}
{\small Il vincolo contribuisce a determinare la decisione e pu\`{o} avere un
proprio valore marginale. }
\end{center}

{\small \smallskip}

{\small \textbf{Lezione 12}: costruiremo il problema duale e useremo i prezzi
ombra per interpretare liquidit\`{a} e funding alle diverse scadenze. }%

%TCIMACRO{\TeXButton{Transition: Box Out}{\transboxout}}%
%BeginExpansion
\transboxout
%EndExpansion%
%TCIMACRO{\TeXButton{EndFrame}{\end{frame}}}%
%BeginExpansion
\end{frame}%
%EndExpansion



%*****************************

\end{document}